% Options for packages loaded elsewhere
\PassOptionsToPackage{unicode}{hyperref}
\PassOptionsToPackage{hyphens}{url}
%
\documentclass[
  10pt,
]{article}
\usepackage{amsmath,amssymb}
\usepackage{iftex}
\ifPDFTeX
  \usepackage[T1]{fontenc}
  \usepackage[utf8]{inputenc}
  \usepackage{textcomp} % provide euro and other symbols
\else % if luatex or xetex
  \usepackage{unicode-math} % this also loads fontspec
  \defaultfontfeatures{Scale=MatchLowercase}
  \defaultfontfeatures[\rmfamily]{Ligatures=TeX,Scale=1}
\fi
\usepackage{lmodern}
\ifPDFTeX\else
  % xetex/luatex font selection
\fi
% Use upquote if available, for straight quotes in verbatim environments
\IfFileExists{upquote.sty}{\usepackage{upquote}}{}
\IfFileExists{microtype.sty}{% use microtype if available
  \usepackage[]{microtype}
  \UseMicrotypeSet[protrusion]{basicmath} % disable protrusion for tt fonts
}{}
\makeatletter
\@ifundefined{KOMAClassName}{% if non-KOMA class
  \IfFileExists{parskip.sty}{%
    \usepackage{parskip}
  }{% else
    \setlength{\parindent}{0pt}
    \setlength{\parskip}{6pt plus 2pt minus 1pt}}
}{% if KOMA class
  \KOMAoptions{parskip=half}}
\makeatother
\usepackage{xcolor}
\usepackage[margin=0.9in]{geometry}
\usepackage{longtable,booktabs,array}
\usepackage{calc} % for calculating minipage widths
% Correct order of tables after \paragraph or \subparagraph
\usepackage{etoolbox}
\makeatletter
\patchcmd\longtable{\par}{\if@noskipsec\mbox{}\fi\par}{}{}
\makeatother
% Allow footnotes in longtable head/foot
\IfFileExists{footnotehyper.sty}{\usepackage{footnotehyper}}{\usepackage{footnote}}
\makesavenoteenv{longtable}
\setlength{\emergencystretch}{3em} % prevent overfull lines
\providecommand{\tightlist}{%
  \setlength{\itemsep}{0pt}\setlength{\parskip}{0pt}}
\setcounter{secnumdepth}{-\maxdimen} % remove section numbering
\usepackage{amsmath,amssymb,amsthm}
\usepackage{longtable}
\ifLuaTeX
  \usepackage{selnolig}  % disable illegal ligatures
\fi
\usepackage{bookmark}
\IfFileExists{xurl.sty}{\usepackage{xurl}}{} % add URL line breaks if available
\urlstyle{same}
\hypersetup{
  pdftitle={Bounded divisorial references and numerical-zero non-Archimedean equations},
  pdfauthor={Research paper with independent scoped review119},
  hidelinks,
  pdfcreator={LaTeX via pandoc}}

\title{Bounded divisorial references and numerical-zero non-Archimedean
equations}
\author{Research paper with independent scoped review119}
\date{5 October 2026}

\begin{document}
\maketitle

\section{Abstract and exact status}\label{abstract-and-exact-status}

Independent review119 returned \textbf{PASS} for two exact extensions:
every bounded genuine change of one prescribed divisorial reference, and
every actual numerically trivial rational bundle on the entire genuine
numerical cone. All indispensable prescribed-reference constructions are
included below. The field is algebraically closed, trivially valued,
characteristic zero and of arbitrary cardinality; the original variety
is smooth integral projective of positive dimension. Raw metric products
and masses are retained.

Original P16.5 remains \textbf{open}. Its general minimal-big and
nonzero numerical pseudoeffective equations are not certified by this
scoped result. The exact actual-ample G0 theorem is the sole equation
grant and is used only within that original ample domain. Historical
author117 candidate=false and original review114 INCONCLUSIVE remain
preserved.

\section{1. Scope, genuine definitions and the allowed ample
input}\label{scope-genuine-definitions-and-the-allowed-ample-input}

Throughout, \(k\) is algebraically closed, trivially valued, of
characteristic zero, with arbitrary cardinality. The original \(X\) is
smooth, integral and projective over \(k\), with dimension \(n\geq1\).
No toric, skeletal, support, countability or finite-generation
hypothesis occurs. Rational bundles are actual elements of
\(\operatorname{Pic}(X)\otimes\mathbb Q\), with their rational
isomorphisms retained; they are not just numerical classes. Write
\(\theta_B=c_1(B)\) in \(N^1(X)\), \(v_0\) for the generic trivial
valuation, and \(X^{\rm an}\) for the compact Hausdorff original
analytic carrier. All measures below are finite Radon measures on that
carrier. Intersections use its fundamental cycle, with no probability
factor.

The canonical trivial metric on an actual bundle is defined by a
norm-one algebraic frame at the valuation-ring center. Properness gives
that center for every analytic point, including nongeneric support.
Units in its local ring have norm one, so the definition is independent
of that center frame. A generating family of sections of a bundle has
\[\max_j\log\|t_j\|_{\rm triv}=0\quad\hbox{on every analytic point}.\tag{1}\]
Indeed all local coefficients have norm at most one and some coefficient
is a unit at the center. Tensor products and rational powers of these
canonical metrics agree literally. The supplied BJ Remark 1.24
additionally gives continuity of the canonical trivial metric on
projective \(X\) and of each section norm into \([0,1]\). These are
metric statements, not assertions of semipositivity for a nonample
bundle.

A real generically finite Fubini--Study function in an actual rational
\(B\) is \[f=\frac1m\max_j\{\log\|s_j\|_{\rm triv}+\lambda_j\},\qquad
0\ne s_j\in H^0(X,mB),\quad \lambda_j\in\mathbb R.\tag{2}\] The section
family need not generate. Values at common zeros are \(-\infty\). On
integral \(X\) every displayed nonzero section is finite at \(v_0\). By
the \textbf{literal single-net definition},
\(\operatorname{PSH}_{\rm def}(\theta)\) consists of generically finite
all-point decreasing limits of one such FS net in actual \(B_i\) with
\(\theta_{B_i}\longrightarrow\theta\). The limit is usc. We write
\(\operatorname{PSH}(L)=\operatorname{PSH}_{\rm def}(\theta_L)\). The
actual metric associated to a scalar \(u\) is
\(\|\cdot\|_u=e^{-u}\|\cdot\|_{\rm triv,L}\). Thus a numerical cone and
an actual metrized bundle remain distinct data even when the former
depends only on \(\theta_L\).

\textbf{Grant G0, used only on actual ample rational bundles.} In the
present field domain, on a reduced irreducible geometrically unibranch
projective carrier with actual ample rational \(H\), the full supplied
G0 theorem provides its genuine fixed-bundle finite continuous FS
approximation; finite maxima and fixed-ample decreasing closure;
all-point divisorial determination and the usc obstacle principle;
compactness of the normalized cone in divisorial topology; envelope
property for continuous obstacles; genuine continuous metric and
finite-energy raw mixed products; arbitrary decreasing-net and
integrable weighted convergence; energy, locality, comparison and
positive-mass domination. Put \(V_H=H^n>0\). Its actual finite-energy
space \(\mathcal E^1(H)\), raw energy \(E_H\), and targets satisfy
\[ E_H(0)=0,\quad E_H(\psi+c)=E_H(\psi)+V_Hc,
\quad \operatorname{MA}_H(\psi)(X^{\rm an})=V_H.\tag{3}\] Targets have
raw mass \(V_H\) and finite dual energy
\[\sup_{\psi\in\mathcal E^1(H)}\left(E_H(\psi)-\int\psi\,d\mu\right)<\infty,\tag{4}\]
equivalently integrability of every genuine \(\mathcal E^1(H)\)
potential. G0 supplies the two-sided outer derivative for every original
continuous direction, bounded-\(J\) coercivity and target-integral
continuity, actual attainment, and the genuine equation unique modulo
constants. Its continuous contact bridge applies to an already glued
continuous rational metric on an actual ample bundle with the specified
finite unit-log same-constant supports. The metric products of canonical
zero ample potentials, on the trivial test configuration, are
\[\bigwedge_{i=1}^n T_{H_i,0}=(H_1\cdots H_n)\delta_{v_0}.\tag{5}\]
Formula (5) is its fundamental-cycle model formula, not a probability
convention. We use the mixed cocycle and integration-by-parts statements
solely with actual ample entries or signed differences of such entries.
No statement of G0 is applied to a nonample polarization merely by
changing its label.

The complete actual-ample grant is the accepted standalone G0 proof
identified in the source provenance. The external primary premises and
their precise applications are itemized in Section 9. Historical
evidence filenames are provenance coordinates, not additional theorem
grants.

\section{2. Direct single-net
foundations}\label{direct-single-net-foundations}

\subsection{2.1. Rational rounding and actual ample
embedding}\label{rational-rounding-and-actual-ample-embedding}

Let \(\phi_i\downarrow\phi\) be one defining real FS net, with degree
\(m_i\) and bundles \(B_i\). For positive rational \(\delta\), round its
coefficients upward to rational numbers with errors at most
\(m_i\delta\). Denote the result by \(f_{i,\delta}\) and set
\(g_{i,\delta}=f_{i,\delta}+\delta\). At every point, including common
poles, \[\phi_i\leq f_{i,\delta}\leq\phi_i+\delta.\tag{6}\] Order pairs
by increasing \(i\), decreasing \(\delta\), and decreasing \(g\)
pointwise. Given two pairs, choose \(i_3\) above both and
\(2\delta_3\leq\min(\delta_1,\delta_2)\). Then
\[g_{i_3,\delta_3}\leq\phi_{i_3}+2\delta_3
\leq\min(\phi_{i_1}+\delta_1,\phi_{i_2}+\delta_2)
\leq\min(g_{i_1,\delta_1},g_{i_2,\delta_2}).\tag{7}\] This defines a
directed decreasing net. Its index map is increasing and final in the
original index set, and \(\delta\to0\). Finite thresholds show
convergence even at \(-\infty\). Thus the class limit is unchanged and
the same all-point endpoint has one rational FS net. There is no strict
comparison of two \(-\infty\) values.

Choose an actual ample rational \(H\) so \(\theta_H-\theta\) is ample.
Eventually \(H-B_i\) is an actual ample rational bundle by openness and
numerical invariance of ampleness. Clear powers and multiply every
displayed section by a generating family of a sufficiently large power
of \(H-B_i\). Identity (1) makes the new FS scalar exactly the old
scalar, including its poles. A tail is therefore a decreasing net in one
actual ample \(H\), proving
\[\operatorname{PSH}_{\rm def}(\theta)\subset\operatorname{PSH}(H).\tag{8}\]
For a finite list of numerical classes one sufficiently large actual
\(H\) works. Hence each defining numerical function is finite on
divisorial valuations, satisfies \(\sup\phi=\phi(v_0)\), and has the
all-point divisorial determination and usc obstacle implication supplied
by G0 in that fixed ample cone. This does not prove decreasing closure
when the numerical classes vary.

\subsection{2.2. A direct finite-maximum closure in a fixed numerical
class}\label{a-direct-finite-maximum-closure-in-a-fixed-numerical-class}

This useful additional closure does not require general flattening.
Suppose \(\phi_i\downarrow\phi\) in \(B_i\), \(\psi_j\downarrow\psi\) in
\(C_j\), with both class limits equal to \(\theta_L\). Fix an actual
ample \(H\). For rational \(\delta>0\), put \(E_\delta=L+\delta H\). On
tails in \(i,j\), both \(E_\delta-B_i\) and \(E_\delta-C_j\) are actual
ample. Multiply the two section families separately by generating
families of these differences, clear to one degree, and take their
union. The resulting actual \(E_\delta\) FS function is exactly
\(\max(\phi_i,\psi_j)\).

The admissible triples \((i,j,\delta)\) are directed by increasing
\(i,j\) and decreasing \(\delta\): choose a smaller \(\delta\) first,
then indices beyond both old indices and its ample-difference tails. The
scalar decreases, converges everywhere to \(\max(\phi,\psi)\), and its
classes tend to \(\theta_L\). Thus finite maxima preserve the literal
genuine cone. Addition of a constant is immediate. Rational convex
combinations follow by tensor products and clearing rational exponents
on the product index set. These direct assertions do not imply closure
of arbitrary infinite suprema or of nested extended-valued varying-class
limits.

\subsection{2.3. Finite-polytope homogeneous
projection}\label{finite-polytope-homogeneous-projection}

For \(\phi(v_0)\geq0\) define
\[\widehat\phi(v)=\inf_{t>0}\frac{\phi(tv)}t.\tag{9}\] It preserves
order, is homogeneous, and commutes with decreasing infima. For rational
FS (2), let \(x_j=\log\|s_j\|(v)\) and put \(r=1/t\). Its finite
computation is \[m\widehat\phi(v)=\inf_{r\geq0}\max_j(x_j+r\lambda_j)
=\max_{w\in\operatorname{Vert} P}\sum_jw_jx_j,
\quad P=\{w\geq0:\sum_jw_j=1,\ \sum_jw_j\lambda_j\geq0\}.\tag{10}\] For
finite \(x_j\), every feasible weighted sum is bounded by the maximum
for every \(r\). A minimum of the convex piecewise affine maximum in
\(r\) occurs at zero, at a breakpoint, or at infinity. At zero an active
nonnegative slope supplies a feasible unit vector. At a finite interior
minimum a zero slope or two opposite active slopes with zero weighted
average supplies an attaining feasible vertex. At infinity a finite
minimum has largest slope zero and an attaining zero-slope unit vector.
These cases prove equality. The vertices are precisely eligible unit
vectors and two-coordinate zero averages of strictly opposite slopes,
hence rational. Feasibility holds because \(\max_j\lambda_j\geq0\).
Continuity at \(r=0\) permits replacing \(r>0\) by \(r\geq0\).

Take a common weight denominator \(b\). Each vertex yields an actual
section \[S_w=\prod_j s_j^{bw_j}\in H^0(X,mbB),\tag{11}\] nonzero on
integral \(X\). Its section logarithm, divided by \(mb\), gives the
weighted expression in (10). Zero exponents mean omission, not
\(0(-\infty)\). When some \(x_j=-\infty\), replace them by
\(\max(x_j,-N)\), commute the infima in \(N,r\), and pass the fixed
finite maximum to its decreasing limit. Thus (10)--(11) hold at all
poles.

Apply this to the rational net of Section 2.1. The homogeneous actual FS
functions decrease to (9), have generic value zero and the original
class limit. Consequently \(\widehat\phi\) belongs to the literal
genuine numerical cone. No theorem about nested nonample limits was
used.

\subsection{2.4. All-point homogeneous decomposition: corrected pole
argument}\label{all-point-homogeneous-decomposition-corrected-pole-argument}

Put \(C=\phi(v_0)\) and \(h_\lambda=\widehat{\phi-\lambda}\) for
\(\lambda<C\). Then
\[\phi(v)=\sup_{\lambda<C}\{h_\lambda(v)+\lambda\}\quad\text{at every point}.\tag{12}\]
The inequality from right to left is immediate from \(t=1\). For the
other direction the ray \(g(t)=\phi(tv)\) is a decreasing limit of
convex decreasing FS rays. Choose one initial tail; all of its rays have
the finite uniform upper bound \(M=\phi_{i_0}(v_0)\).

If \(g(t_0)=-\infty\), monotonicity handles \(s>t_0\). For \(0<s<t_0\)
choose a \textbf{still smaller} auxiliary parameter \(0<a<s\). Write
\(s=(1-\beta)a+\beta t_0\) with \(0<\beta<1\). At each stage,
\[g_i(s)\leq(1-\beta)g_i(a)+\beta g_i(t_0)
\leq(1-\beta)M+\beta g_i(t_0).\tag{13}\] The right side tends to
\(-\infty\). Hence the entire ray is \(-\infty\), and (12) follows
there. This is the requested correction to the earlier erroneous word
``larger.'' It does not interpolate an undefined expression at the
limit.

Otherwise the ray is finite convex. A supporting line \(at+\lambda_0\)
at \(t=1\) has \(a\leq0\) and \(\lambda_0\leq C\), since \(g(t)\leq C\).
If \(\lambda_0<C\), definition (9) gives \(h_{\lambda_0}(v)\geq a\),
proving equality in (12). If \(\lambda_0=C\), lower the intercept to
\(C-\varepsilon\), keep the same slope, and obtain a right side at least
\(g(1)-\varepsilon\). Let \(\varepsilon\downarrow0\). This proves (12)
without restricting the support of \(v\).

\section{3. Reconstructed all-big prescribed-divisor
theorem}\label{reconstructed-all-big-prescribed-divisor-theorem}

\subsection{3.1. Actual data for every big rational
bundle}\label{actual-data-for-every-big-rational-bundle}

Fix an actual rational \(L\) with big numerical class. Choose an actual
ample integral \(H\). For sufficiently small rational \(\varepsilon>0\),
\(\theta_L-\varepsilon\theta_H\) is pseudoeffective. Approximate this
class by finite nonnegative rational combinations of effective Cartier
divisors closely enough that \(\theta_L-\theta_D\) is ample. Such
approximation follows from the definition of the closed effective cone,
rational approximation of each finite coefficient list, and
finite-dimensional openness of ampleness. Smooth local factoriality
ensures that the components used are Cartier. Set the \textbf{actual}
bundle \(A=L-\mathcal O(D)\). Numerical ampleness makes it actual ample.
Thus, on the same original \(X\),
\[L\simeq A+\mathcal O(D),\qquad A\text{ actual ample rational},\quad D\geq0\text{ rational Cartier}.\tag{14}\]
This treats numerically trivial twists of big actual bundles as well. No
Bertini theorem, modification, section-ring finite generation, or field
cardinality restriction is introduced. Conversely an actual ample bundle
plus an effective rational divisor is big: small rational perturbations
of an ample class remain ample and have sections after a power, giving
an open neighborhood in the effective cone. Adding an effective class
preserves its interior.

Fix (14) \textbf{before} a target is selected. Put
\(\sigma=\log\|s_D\|_{\rm triv}\), \(Z=|D|^{\rm an}\),
\(U=X^{\rm an}\setminus Z\), and \(V_D=A^n>0\). If \(D=0\), \(\sigma=0\)
and \(Z\) is empty. Define the whole genuine reference-dominated cone
\[\mathcal P_D(L)=\{u\in\operatorname{PSH}(L):u\leq\sigma+C\text{ everywhere for some }C\in\mathbb R\}.\tag{15}\]

\subsection{3.2. Entire cone
reconstruction}\label{entire-cone-reconstruction}

For every actual \(A\)-psh \(\psi\), multiply its decreasing finite
fixed-\(A\) FS sections by a cleared canonical section of \(D\). This
gives one decreasing actual \(L\) FS net with endpoint \(\sigma+\psi\),
including the value \(-\infty\) on \(Z\). Thus \(F_D(\psi)=\sigma+\psi\)
belongs to (15).

For the converse first let \(h\) be homogeneous genuine \(L\)-psh with
generic value zero and \(h\leq\sigma\). Suppose \(D\ne0\). Section 2.3
gives a decreasing homogeneous rational FS net \(h_i\) in \(B_i\) with
classes tending to \(\theta_L\). The set \(K=\{\sigma=-1\}\) is compact
and nonempty. For rational \(0<t<1\), the open sets \(\{h_i<-t\}\cap K\)
increase and cover \(K\), so an entire tail satisfies that strict
inequality. On the cofinal admissible pairs \((i,t)\), with \(i\)
increasing and \(t\uparrow1\), homogeneity gives \(h_i\leq t\sigma\) on
divisorial points, by scaling to \(K\) when \(\sigma<0\) and using
\(h_i\leq0\) when \(\sigma=0\). The ample embedding and usc obstacle
principle extend this inequality to every point.

Clear rational powers. Each displayed section in \(h_i\) vanishes to the
required order along each prime of \(tD\). The normal section-division
fact, whose complete supplied proof is BJ Lemma 2.23, gives an actual
homogeneous quotient FS function \(r_{i,t}\) in \(B_i-tD\), with generic
value zero, satisfying \(r_{i,t}=h_i-t\sigma\) on \(U\). It is defined
by divided sections on all of \(X^{\rm an}\), not by subtracting two
infinite values. \textbf{The divided net is not asserted to decrease.}

Adjoin rational \(\eta>0\), put \(c=1+\eta\), and require
\(cA-(B_i-tD)\) actual ample. These admissible triples are directed and
cofinal with \(i\) increasing, \(t\uparrow1\), \(\eta\downarrow0\):
choose a smaller \(\eta\), then \(t\) close enough to one and \(i\)
beyond its numerical ample tail. Multiplication by generators and
rational scaling make \(r_{i,t}/c\) actual normalized \(A\)-FS
functions. G0 compactness supplies a subnet with actual \(A\)-psh limit
\(r\). On each divisorial point its limit is explicitly \(h-\sigma\),
since both values are finite. The forward construction makes
\(\sigma+r\) genuine \(L\)-psh. Place this and \(h\) in one actual ample
\(H\) by (8). Divisorial determination gives
\[h=\sigma+r\quad\text{everywhere}.\tag{16}\] The same determination in
\(A\) gives uniqueness of \(r\). When \(D=0\) the assertion is simply
the actual ample identification.

Now let \(u\) satisfy (15). Its homogeneous decomposition (12) has
\(h_\lambda+\lambda\leq u\leq\sigma+C_0\). Evaluate at \(tv\) for a
divisorial \(v\), divide by \(t\) and let \(t\to\infty\). This gives
\(h_\lambda\leq\sigma\) divisorially, then everywhere by (8) and the usc
obstacle principle. By (16), \(h_\lambda=\sigma+r_\lambda\) with
normalized actual \(A\)-psh \(r_\lambda\). The functions
\(r_\lambda+\lambda\) are uniformly bounded above by \(C=u(v_0)\). Their
finite maxima \(b_J\) are actual \(A\)-psh with generic value
\(c_J=\max J\uparrow C\). Compactness of \(b_J-c_J\) gives an actual
\(A\)-psh \(r\), with divisorial values
\(\sup_{\lambda<C}(r_\lambda+\lambda)\) and generic value \(C\). Each
constituent is at most \(r\) by all-point divisorial domination.
Conversely \(r\) is at most the usc regularization of their supremum by
the usc obstacle principle. Thus \(r\) is that genuine ample supremum.
On divisorial points, (12) gives \(u=\sigma+r\); ample determination in
one common \(H\) extends this equality everywhere.

Consequently \[F_D:\operatorname{PSH}(A)\longrightarrow\mathcal P_D(L)
\quad\text{is an order bijection}.\tag{17}\] Its inverse \(Q_D\) is the
unique actual \(A\)-psh extension of \(u-\sigma\) from \(U\). For
injectivity or order reflection subtract the finite \(\sigma\) only on
divisorial points and use G0. In particular the quotient has uniquely
determined values on \(Z\); no arbitrary lift or modification appears.

\subsection{3.3. Compactness and envelopes, including the empty
case}\label{compactness-and-envelopes-including-the-empty-case}

One has \(\sup u=u(v_0)=Q_Du(v_0)=\sup Q_Du\). On each divisorial
coordinate the map (17) is addition of a fixed finite scalar; hence it
and its inverse are homeomorphisms. G0 normalized compactness proves
compactness of the actual original normalized cone. If a decreasing net
in this fixed reference cone has a generically finite endpoint, order
reflection and ample decreasing closure identify its quotient limit
everywhere. This is fixed-reference closure only.

For a bounded-above usc obstacle \(b\), feasibility is equivalent in the
two cones, since a divisorial inequality \(Q_Du\leq b\) implies it
everywhere. If there are no minorants, the envelope is \textbf{by
convention \(-\infty\)}, outside the generically finite cone. With that
convention, \[P_D(\sigma+b)=\sigma+P_A(b).\tag{18}\] For nonempty
uniformly bounded families, usc suprema commute with \(F_D\): addition
is finite continuous on \(U\), and the common upper bound on the
quotients forces value \(-\infty\) at \(Z\) as \(\sigma\to-\infty\).
Continuous finite \(b\) have constant ample minorants, so (18) has a
continuous genuine quotient by G0 EP. This is the requested harmless
empty-feasible-cone clarification.

\section{4. Genuine original metric products and all
limits}\label{genuine-original-metric-products-and-all-limits}

\subsection{4.1. Local metric identity and independent generated
approximations}\label{local-metric-identity-and-independent-generated-approximations}

Clear powers and write \(s_D=f e_D\) and \(e_L=e_Ae_D\) in actual
frames. The local logarithmic metric weights of \(u=\sigma+\psi\)
satisfy \[h_L=h_A+\log|f|.\tag{19}\] The sign is plus. On \(U\) the
nowhere vanishing section \(s_D\) identifies the actual metrized \(A\)
with \(L\): the frame \(e_As_D=f e_L\) has weight exactly \(h_A\). CLD's
supplied harmonic unit-log statement and continuous product invariance
therefore identify the actual continuous products on \(U\). BE's
supplied semiample smooth regularization and continuous product theorem
apply to the actual ample \(A\), including trivial valuation. They
identify these products with the raw G0 products before any target
measure is selected.

It is necessary to justify independence of \textbf{independently
supplied local} bounded decreasing generated approximations, not merely
global approximants. We give the common-collar argument on the original
projective \(X\), with the identity modification. The category consists
of finite rational affine FS/contact weights on a common ordinary
neighborhood, with bounded decreasing endpoints. It includes the local
section weights of our metrics, their bounded quotient floors, and their
finite continuous generated references. It is not a density claim for
arbitrary local positive currents.

Near a tested point of \(U\), shrink to a principal algebraic chart in
which the frames in (19) and \(f,f^{-1}\) are defined. Finitely many
strict rational bounds on its regular coordinates specify a smaller
ordinary neighborhood. A principal chart of a very ample system writes
the used regular functions as \(F_j/a^m\) after clearing denominators.
Append a complete sufficiently ample section system; its ratios generate
that chart ring. Negative rational constant floor and strict margins
give one actual ample \(H\) and a finite global FS potential \(p\),
whose local raw reference \(v=h_0+p\) has
\[Q=\{v\leq0\}\Subset U_{\rm chart},\quad
P=\{v\leq-\epsilon\},\quad x\in\operatorname{int}P,\quad\epsilon>0.\tag{20}\]
Explicitly, with cleared tensor denominator \(q\),
\[v=\frac1q\max(q\lambda,\log|F_j/a^m|-C_j).\tag{21}\] The thresholds
are \(C_j\), not \(qC_j\). The same expression is obtained from actual
sections in the \(a^m\) frame. The appended generator bounds make \(Q\)
compact; principal denominators and strict margins ensure it lies in the
chosen ordinary neighborhood. Bounds on inverses are included by
shrinking the algebraic chart first. This is precisely the actual ample
collar construction in the supplied G0 body, applied only to an actual
ample bundle on projective \(X\). The local \(h_0\) is finite continuous
on \(Q\).

Let \(w_i\) be a finite tuple of bounded raw local weights on a
neighborhood of \(Q\), supplied as decreasing finite rational affine
FS/contact nets \(w_{i\alpha}\). Include two independent nets whenever
one endpoint is to be compared, and include a bounded generated test if
needed. Choose rational \(M>0\) with \(-M\leq w_i\leq M\) on \(Q\), and
rational \(R\) with \(R\epsilon>2M+2\). Set \(h=Rv+M+1\). On \(P\) it is
below \(-M\). On the compact outer strip
\[S=Q\cap\{v\geq-1/(2R)\},\qquad h\geq M+1/2.\tag{22}\] For each net the
closed sets \(S\cap\{w_{i\alpha}\geq M+1/4\}\) decrease to the empty
set. Compactness gives an entire tail empty on the whole strip. This
argument does not assert uniform convergence to a discontinuous
endpoint. Choose these tails for the finite tuple separately.

Define global functions
\[\Psi_{i\alpha}=\max(w_{i\alpha}-h,0)\text{ on }Q,
\quad\Psi_{i\alpha}=0\text{ outside }Q,
\quad\Phi_{i\alpha}=p+R^{-1}\Psi_{i\alpha}.\tag{23}\] They are
continuous and glued on the same actual \(H\): a whole fixed outer
collar has \(\Psi=0\), and the inner raw weight is
\[R^{-1}\max(w_{i\alpha}-M-1,Rv).\tag{24}\] Deleting inactive
zero-coordinate branches on smaller unit charts, these weights have the
exact same-constant affine supports required by G0's continuous contact
bridge. Thus \(\Phi_{i\alpha}\) is genuinely \(H\)-psh. It decreases to
the explicit bounded global endpoint \(\Phi_i=p+R^{-1}\max(w_i-h,0)\) on
\(Q\), and \(p\) outside. G0 fixed-ample closure makes that endpoint
genuine finite energy. On \(\operatorname{int}P\) one has, at finite
stages and at the endpoint, \[w_i=R\Phi_i+Rh_0+M+1,\qquad
dd^cw_i=R(c_1(H)+dd^c\Phi_i).\tag{25}\] The fixed reference term
\(Rh_0\) is indispensable.

The finite-stage continuous products are the actual original products by
BE/CLD. For compactly supported continuous tests in the inner open,
extend by zero and use G0 mixed decreasing-net convergence for the
\(\Phi_{i\alpha}\). The local limit is \(R^n\) times the restriction of
the genuine raw \(H\) mixed measure. Independent nets with the same
endpoints use the \textbf{same} \(p,h,H,R\) and therefore have exactly
the same global endpoints in (23). Their limiting products agree. Two
different collars are compared on a third smaller shared collar, where
their finite-stage products agree; compact tests pass this to their
bounded limits. Positivity and the local Radon property follow from the
ample products, so these limits glue.

For a bounded generated test \(w_0\), (25) splits its weighted pairing
into \(R\Phi_0\) and a fixed continuous local reference term. G0's
weighted convergence applies to the first part with the sign
\[\text{test-containing intersection pairing};-;\text{zero-test pairing}.\tag{26}\]
Ordinary weak convergence applies to the reference part. Signed
differences are handled separately; a difference of decreasing nets is
not called decreasing.

Real constants do not change this conclusion. Round up with error
\(\delta\) and add \(\delta\) as in (6)--(7). For each fixed real stage
the explicit global formulas (23) are approached uniformly by rational
stages, with error at most \(2\delta/R\), while preserving the collar
margin. Fixed-ample uniform closure and continuous product continuity
make the real stage genuine. The original real-stage net then decreases
to exactly the same global endpoint, so G0 applies directly, without
interchanging unrelated limits.

On a frame overlap, transform a whole chosen generated net by the same
harmonic \(\log|g|\) for the actual invertible transition function
\(g\). Its endpoint is the transformed weight. Finite-stage products
agree by CLD; compact tests pass the equality to the local limit.
Independence from other generated nets removes the choice of the matched
net. This establishes actual frame compatibility before the unbounded
construction.

One collar commutes with maxima, since
\(\Phi_{\max(w,z)}=\max(\Phi_w,\Phi_z)\). G0 strict-winning locality
therefore gives local locality on the actual Borel winning sets, even
for discontinuous bounded endpoints. For different finite continuous
generated references \(b,b'\), bound \(|b-b'|\leq C\) on \(Q\), and
compare floors using
\[z_\ell=\max(w,b-\ell,b'-\ell),\qquad\ell>\max(j,k)+C.\tag{27}\] Then
\(\max(w,b-j)=\max(z_\ell,b-j)\) and similarly for \(b'-k\).
Strict-winning locality compares their products with that of \(z_\ell\)
on the respective winning sets. Those sets exhaust \(\{w>-\infty\}\),
and later restricted products stabilize on every fixed winning set.
Their increasing nonpluripolar local limits agree. The same construction
for a finite tuple gives mixed compatibility. No unrelated classical
local-current density theorem is used.

\subsection{4.2. An explicit finite-energy polar
test}\label{an-explicit-finite-energy-polar-test}

We must show that the reference divisor and the relevant quotient poles
have zero mass, rather than assume it. In any actual ample \(A\), choose
a high \(q\) and a nonzero section \(s\) of \(qA\) vanishing on \(|D|\).
The supplied ample-twist global generation applied to its coherent ideal
gives such a section, not identically zero on the divisor complement;
additional zeros are allowed. Put \(a=q^{-1}\log\|s\|\), with generic
value zero. The floors \(\max(a,-j)\) are finite actual \(A\)-FS
functions, by adjoining a generating family with constant weight \(-j\)
after clearing powers.

Use normalized probability measures \textbf{only} in this positive-mass
ample estimate, denoted
\(\operatorname{MA}_{\rm BJ,A}=A^{-n}\operatorname{MA}_A\). Set
\[\operatorname{Cap}_A(B)=\sup_{-1\leq b\leq0}\operatorname{MA}_{\rm BJ,A}(b)(B).\tag{28}\]
G0 ample integration by parts and telescoping from the zero metric give
\[\int(-a)\operatorname{MA}_{\rm BJ,A}(b)\leq n,
\qquad \operatorname{Cap}_A\{a\leq-t\}\leq n/t.\tag{29}\] For the first
bound, start with a bounded floor of \(a\). The zero metric measure is
\(\delta_{v_0}\), whose integral of \(-a\) is zero. Each of the \(n\)
telescoping mixed differences is, after integration by parts, the
difference of integrals of \(-b\in[0,1]\) against two positive mixed
measures of normalized mass one, hence has absolute value at most one.
Pass floors by monotone convergence and G0 weighted continuity. The
second bound is then Markov's inequality uniformly over \(b\).

For normalized ample \(c\leq0\), put \(c_s=\max(c,-s)\),
\(\mu_s=\operatorname{MA}_{\rm BJ,A}(c_s)\). Ample energy comparison
gives \(J(c_s)\leq\int(-c_s)d\mu_s\). Locality and equal total mass give
\(\mu_s\{c\leq-t\}=\mu_t\{c\leq-t\}\) when \(1\leq t\leq s\). The
positive mixed expansion of \(\operatorname{MA}_{\rm BJ,A}(t^{-1}c_t)\)
contains \(t^{-n}\mu_t\). Layer-cake integration therefore yields
\[J(c)\leq1+\int_1^\infty t^n\operatorname{Cap}_A\{c\leq-t\}\,dt.\tag{30}\]
The capacity weight is \(t^n\); it must not be dropped.

Choose \(0<\alpha<1/(n+1)\) and \(\chi(x)=1-(1-x)^\alpha\) for
\(x\leq0\). It is convex increasing with slopes in \([0,1]\). Decreasing
piecewise affine chord majorants on expanding intervals, refined at each
step and with constant left tails, give decreasing genuine ample
approximations to \(c=\chi(a)\): each affine branch of slope in
\([0,1]\) is psh by convexity with zero, and finite maxima remain psh.
Thus G0 fixed-ample closure applies. Since
\[\operatorname{Cap}_A\{\chi(a)\leq-t\}=O(t^{-1/\alpha}),\tag{31}\] (30)
proves \(\chi(a)\in\mathcal E^1(A)\), and its pole set contains \(Z\).
For a quotient's own pole set use its normalized ample potential in
place of \(a\) and the same estimates. G0 mixed finite-energy
integrability now implies that all finite-energy mixed ample products,
and every target integrating all ample \(\mathcal E^1\), give these pole
sets zero mass. No countable extraction is involved.

\subsection{4.3. Reference floors, raw mass and unbounded
limits}\label{reference-floors-raw-mass-and-unbounded-limits}

For \(u=F_D(\psi)\) define \(\psi_j=\max(\psi,-j)\),
\(u_j=\max(u,\sigma-j)=F_D(\psi_j)\), and \(B_j=U\cap\{\psi>-j\}\). The
bounded local products just constructed and strict-winning locality give
increasing measures \(\mathbf1_{B_j}T_{u_j}^n\), each with mass at most
\(A^n\). Their increasing finite limit is Radon: on a Borel set choose a
late measure approximating its limiting mass, then approximate that
measure from inside by a compact set; finite total mass and complements
give outer regularity. Define the genuine reference nonpluripolar
operator by this local construction,
\[\operatorname{MA}_D(u)=\lim_j\mathbf1_{B_j}T_{u_j}^n.\tag{32}\] It was
constructed before a target. For \(\psi\in\mathcal E^1(A)\), locality
and the polar test show that the union of the \(B_j\) has full ample
endpoint measure. Thus
\[\operatorname{MA}_D(F_D\psi)=\operatorname{MA}_A(\psi),
\qquad\operatorname{MA}_D(F_D\psi)(X^{\rm an})=A^n.\tag{33}\] This
equality is a derived local metric identity and pole removal, not an
inverse operator definition.

For finitely many data \(L_i=A_i+D_i\), use the common complement of
their divisors and simultaneous quotient floors. Mixed strict-winning
locality follows by finite algebraic polarization, while positivity
comes from the actual ample mixed products, not the alternating
polarization coefficients. The limit for finite-energy quotients is
their genuine ample mixed product, of raw mass \((A_1\cdots A_n)\),
extended on the original \(X\) with zero mass on all removed poles.

For decreasing relative nets with finite-energy endpoints, (17) and
fixed-ample decreasing closure give decreasing quotient nets on every
point. G0 arbitrary unbounded mixed convergence therefore supplies the
original mixed weak limits. For a decreasing integrable ample quotient
test \(\psi_0\), its weighted integral is
\[\int\psi_0\,\bigwedge_{i=1}^n T_{u_i}
=(\theta_{A_0},\psi_0)\prod_i(\theta_{A_i},\psi_i)
 -(\theta_{A_0},0)\prod_i(\theta_{A_i},\psi_i).\tag{34}\] G0 integrable
weighted decreasing continuity gives weak convergence with this
subtraction order. The local raw test differs from its quotient by a
fixed finite continuous reference term, which passes by ordinary weak
convergence on compact frame opens. No assertion of global absolute
integrability of \(\sigma\) is made or required.

\section{5. Energy, targets, all continuous derivatives and attained
solutions}\label{energy-targets-all-continuous-derivatives-and-attained-solutions}

For bounded quotients the genuine products constructed in Section 4
define
\[E_D(u)=\frac1{n+1}\sum_{j=0}^n\int Q_Du\;T_u^j\wedge T_\sigma^{n-j}.\tag{35}\]
Each integral is on the original carrier, using the uniquely retained
quotient. Identities (19), (33) and the mixed version give
\(E_D(u)=E_A(Q_Du)\). For arbitrary \(u\) take decreasing reference
floors in (35); G0 monotonicity and decreasing energy continuity give
\[E_D(u)=E_A(Q_Du),\quad
\mathcal E_D^1=F_D(\mathcal E^1(A)),\quad E_D(\sigma)=0,
\quad E_D(u+c)=E_D(u)+A^n c.\tag{36}\] This equality follows after
actual products have been constructed. Concavity, monotonicity, usc and
mixed integrability are inherited by the proved order homeomorphism,
with the explicit raw mass.

The intrinsic targets for this retained reference have
\[\mu(X^{\rm an})=V_D=A^n,\qquad
D_D(\mu)=\sup_{u\in\mathcal E_D^1}\left(E_D(u)-\int Q_Du\,d\mu\right)<\infty.\tag{37}\]
The condition uses no chosen solution and is equivalent by (17), (36)
and G0 to integrability of every actual \(\mathcal E_D^1\) quotient. It
implies \(\mu(Z)=0\) by Section 4.2. It imposes no separate support
condition and does not require absolute integrability of the original
singular reference. The target mass is compulsory: if it differs from
\(V_D\), translations make the variational supremum infinite. Dividing
by \(V_D\) is permissible in the following ample estimates because
\(V_D>0\), never in a zero-volume case.

For \(u\in\mathcal E_D^1\), \(\psi=Q_Du\) and \(f\in C^0(X^{\rm an})\),
continuous original majorants \(\sigma+b\geq u+tf\) are equivalent to
\(b\geq\psi+tf\). The forward implication is obtained on divisorial
points and then by the ample usc obstacle principle applied to
\(\psi\leq b-tf\). Thus the retained-reference outer energy is
\[E_D^\downarrow(u+tf)=E_A^\downarrow(\psi+tf),\qquad
\left.\frac d{dt}\right|_{0}E_D^\downarrow(u+tf)
=\int f\,\operatorname{MA}_D(u).\tag{38}\] All continuous directions on
the original carrier and both derivative signs are supplied by G0 on the
actual ample \(A\). Values hidden by \(-\infty\) on \(Z\) are retained
by the unique quotient, not assigned afresh. For uniform PL
approximations \(f_m\) the energy error is at most
\(V_D|t|\|f-f_m\|_\infty\) by monotonicity and translation. Dividing for
both signs and passing \(m\) confirms that no continuous direction or
derivative sign is lost.

For completeness actual attainment is a separate step. Let a maximizing
sequence for (37) have normalized quotients \(\sup\psi_j=0\), let
\(\bar\mu=\mu/V_D\), and put \(d=D_D(\mu)/V_D\). The exact ample G0
coercivity gives \[J(\psi_j)\leq C_n\{J_{\bar\mu}(\psi_j)+d\},\quad
J_{\bar\mu}(\psi)=d-V_D^{-1}E_A(\psi)+\int\psi\,d\bar\mu.\tag{39}\] The
second expression tends to zero along this sequence. The \(J\)'s are
bounded. Compactness of the actual normalized cone supplies a subnet
with actual quotient limit \(\psi\). Lower semicontinuity of \(J\) makes
it finite energy. G0 target-integral continuity on bounded-\(J\) sets
and energy upper semicontinuity show that \(u=F_D(\psi)\) attains (37).
A sequence was chosen only to approximate a real supremum; no sequential
compactness or countability of \(k\) is assumed.

Let \(P_D(u+tf)\) denote the relative envelope, whose quotient is
\(P_A(\psi+tf)\). It is finite energy, bounded in difference from
\(\psi\), and is at most its obstacle. Its objective is at least the
outer objective, hence maximality implies
\[E_D^\downarrow(u+tf)-\int(\psi+tf)d\mu
\leq E_D(u)-\int\psi d\mu.\tag{40}\] For positive \(t\), (38)--(40) give
\(\int f\operatorname{MA}_D(u)\leq\int f d\mu\); negative \(t\) gives
the reverse inequality. Equality for every continuous \(f\) identifies
the already constructed original operator with \(\mu\) by uniqueness of
Radon measures. Conversely ample concavity proves every finite-energy
solution attains this objective and has finite dual energy. The
intrinsic targets are therefore exactly the solved measure class.

On \(U\), strict inequalities of \(u,v\) are precisely those of their
quotients, and their measures give \(Z\) zero mass. Ample comparison
yields \[\int_{\{u<v\}}\operatorname{MA}_D(v)\leq
\int_{\{u<v\}}\operatorname{MA}_D(u).\tag{41}\] Positive-mass domination
follows from quotient domination and order reflection. Here \(I_A\)
denotes the probability-normalized ample functional, with measure
\(V_D^{-1}\operatorname{MA}_A\); the original relative functional
retains raw mass. The actual relative \(I\) functional is
\[I_D(u,v)=\int(Q_Du-Q_Dv)(\operatorname{MA}_D(v)-\operatorname{MA}_D(u))
=V_D I_A(Q_Du,Q_Dv)\geq0.\tag{42}\] It never subtracts infinite original
scalars. Equal measures force equal ample quotients modulo one real
constant, hence all-point equality \(u=v+c\). This proves the entire
prescribed-cone equation in all big actual rational bundles, with raw
reference mass \(A^n\); it does not identify that mass with
\(\operatorname{vol}(L)\).

\section{6. New bounded-reference compatibility
theorem}\label{new-bounded-reference-compatibility-theorem}

\textbf{Theorem 6.1.} Retain any fixed datum (14). Let
\(\tau\in\operatorname{PSH}(L)\) have the same bounded singularity type
as \(\sigma\), meaning
\[\sigma-C\leq\tau\leq\sigma+C\quad\text{everywhere}\tag{43}\] for
finite \(C\), with inequalities interpreted at poles. Then the entire
genuine dominated cone, its finite-energy space, genuine nonpluripolar
operator, target mass and finite-energy target condition agree for the
two references. The relative energies have the explicit cocycle below.
This applies to every such genuine \(\tau\), not only to a continuous
section reference.

By (17), \(\tau=F_D(\chi)\) for a unique actual ample \(\chi=Q_D\tau\).
On divisorial points (43) gives \(|\chi|\leq C\). G0 obstacle comparison
extends both bounds everywhere, so \(\chi\) is bounded. Conversely every
bounded ample \(\chi\) gives a reference (43). Dominating by
\(\tau+O(1)\) is equivalent to dominating by \(\sigma+O(1)\), so their
genuine cones coincide. The retained relative difference for
\(u=F_D(\psi)\) is the uniquely defined \(\psi-\chi\), including on
\(Z\), not the indeterminate original subtraction \(u-\tau\) there.

Use floors \(\max(u,\tau-j)=F_D(\max(\psi,\chi-j))\). Their bounded
original products are already independently constructed by Section 4.
The sets \[U\cap\{\psi>\chi-j\}\] increase to the finite quotient locus.
For finite-energy \(\psi\), ample strict-winning locality identifies the
restricted floor product with \(\operatorname{MA}_A(\psi)\) on these
sets. Polar removal gives full mass, so the reference-\(\tau\) operator
equals (33).

Independence holds for all endpoints of the associated nonpluripolar
construction, not just full-mass ones. Compare two bounded references
\(\chi,\chi'\) using the common bounded lower floor
\(\max(\psi,\chi-\ell,\chi'-\ell)\), with
\(\ell>j,k+\|\chi-\chi'\|_\infty\). G0 Borel strict-winning locality
makes the restricted products agree on their overlapping winning sets.
These sets exhaust \(\{\psi>-\infty\}\); positivity and increasing
limits give the same quotient nonpluripolar limit. On the reference
divisor both constructions extend with zero mass. Mixed products use the
joint winning set of a finite tuple. Thus the operator compatibility is
obtained from actual local products and floor comparisons, without
assuming an unrestricted local-current converse.

The actual relative mixed energy based at \(\tau\) is, first for bounded
\(\psi\) and then by these floors, \[E_\tau(u)=\frac1{n+1}\sum_{j=0}^n
\int(\psi-\chi)\;T_u^j\wedge T_\tau^{n-j}
=E_A(\psi)-E_A(\chi)=E_D(u)-E_D(\tau).\tag{44}\] The middle equality is
the actual ample cocycle, applied only in \(A\). All terms for
finite-energy \(\psi\) are integrable by G0; bounded \(\chi\) and (34)
justify the floor passage. Therefore
\(\mathcal E^1_\tau=\mathcal E^1_D\) and \(E_\tau(\tau)=0\). For every
target of mass \(V_D\),
\[D_\tau(\mu)=D_D(\mu)-E_D(\tau)+\int\chi\,d\mu.\tag{45}\] The last term
is finite since \(\chi\) is bounded. Thus finite dual energy and
all-finite-energy relative-difference integrability are equivalent for
both references. Formula (45) also shows that the maximizing actual
functions are the same. Their equation, all continuous outer
derivatives, comparison and uniqueness follow from the same original
operator: outer energies differ by \(E_D(\tau)\), a constant independent
of \(t\).

This proves compatibility throughout the bounded singularity type of the
prescribed divisor. It does not say that all references in a big
numerical class are boundedly equivalent or share one full mass.

\subsection{6.2. A reference-mass family on every big actual
bundle}\label{a-reference-mass-family-on-every-big-actual-bundle}

There is a precise reason not to infer a class-independent full mass
from Theorem 6.1. Start with any (14), and choose a nonzero section of a
high integral \(qA\), whose divisor is \(E\geq0\). Positive dimension
and \(A^n>0\) force \(E\ne0\): a nowhere vanishing section would
trivialize \(qA\). For any rational \(0<t\leq1\), the actual data
\[A_t=tA,\qquad D_t=D+(1-t)E/q,\qquad L\simeq A_t+D_t\tag{46}\] have
reference mass \(t^n A^n\). The rational isomorphism uses the chosen
actual section of \(qA\), not only its numerical class. The references
for \(t<1\) and \(t=1\) are not boundedly equivalent: a divisorial
valuation with positive value on \(E\) and its scaling ray make
\(\sigma_{D_t}-\sigma_D\) unbounded below. This construction applies
also when the original big \(L\) is nonample. At \(t=0\) the
full-divisor cone below has zero mass, even if this actual \(L\) is big.

These varying masses are consistent with one genuine nonpluripolar
operator having different masses on different singularity types. They
obstruct an assertion that all these reference spaces are a single
full-mass space of mass \(\operatorname{vol}(L)\), or that their dual
target conditions may be identified without proof. No nonexistence of a
more general class operator is inferred.

\section{7. Reconstructed effective full-divisor zero-mass
theorem}\label{reconstructed-effective-full-divisor-zero-mass-theorem}

Let \(D\geq0\) be any effective rational Cartier divisor and retain the
actual \(L=\mathcal O(D)\) and \(\sigma=\log\|s_D\|\). Define the whole
literal dominated cone as in (15), without declaring it a copy of
constants. We prove
\[\mathcal P_D(\mathcal O(D))=\{\sigma+C:C\in\mathbb R\}.\tag{47}\] For
a homogeneous \(h\leq\sigma\), use the homogeneous rational FS net,
compact \(K\), admissible \((i,t)\) and actual section division from
Section 3.2. The quotients \(r_{i,t}\) have classes
\(\theta_{B_i-tD}\to0\) and generic values zero, with finite divisorial
values \(h_i-t\sigma\). Fix an actual ample \(A_0\), and refine by
rational \(\epsilon>0\) tending to zero, requiring
\(\epsilon A_0-(B_i-tD)\) actual ample. This is possible on tails for
each \(\epsilon\). The triples are directed: choose a smaller
\(\epsilon\) first, then sufficiently late \(i,t\) beyond the old pairs
and its ample tail. Multiplication by generators makes
\(r_{i,t}/\epsilon\) normalized actual \(A_0\)-FS.

Evaluation at any fixed divisorial \(v\) is a finite continuous real
function on the compact normalized ample cone. It therefore has a finite
lower bound \(M_v\). Thus \[\epsilon M_v\leq r_{i,t}(v)\leq0.\tag{48}\]
On the refined net this yields \(h(v)=\sigma(v)\). Both genuine
functions embed into one ample \(H\), so equality holds at every point.
For \(D=0\), the direct numerical-zero proof below supplies the same
conclusion without \(K\) or division. For arbitrary \(u\leq\sigma+C_0\),
(12) and scaling give \(h_\lambda\leq\sigma\), hence
\(h_\lambda=\sigma\). Taking (12) gives \(u=\sigma+u(v_0)\) everywhere.
The converse is one actual divisor-section FS function plus a constant.
This proves (47), including pole values.

On the complement of \(D\), the actual raw metric weight of \(\sigma+C\)
in a frame \(s_D=f e_L\) is \(\log|f|+C\). It has zero Hessian by the
supplied CLD analytic-unit identity. Its genuine positive-degree local
products are zero. Reference floors have the same form with \(C\)
replaced by \(\max(C,-j)\). Their strict-winning nonpluripolar extension
is consequently the zero original Radon measure, constructed before any
target. Actual relative mixed energy is zero, the finite-energy space is
the entire cone (47), and the relative quotient is the finite constant
\(C\). All mixed, decreasing-net and integrable weighted limits are
zero.

The normalized cone is a singleton. For continuous \(b\), the exact
envelope is \(\sigma+\inf_X b\): feasibility forces \(C\leq b\) on
divisorial points and then everywhere by the ample usc obstacle
principle. For an arbitrary usc obstacle the empty case has envelope
\(-\infty\). For \(u=\sigma+C\) and every continuous original \(f\), the
minorant \(\sigma+C-|t|\|f\|_\infty\) makes the feasible cone nonempty,
so its outer energy is zero and its two-sided derivative is
\(0=\int f\,0\). The sole positive target of mass zero is \(\mu=0\),
with zero intrinsic dual energy and automatic quotient integrability.
Every \(\sigma+C\) is an actual attained solution, unique modulo
constants within this whole prescribed cone. Comparison and locality are
identities of zero measures. \textbf{Positive-mass domination is not
claimed:} its almost-everywhere premise would be vacuous at zero mass
and cannot imply an order inequality.

The construction (47) works even when \(\mathcal O(D)\) is big, but its
stated additional \textbf{zero-volume} domain for P16.5 is restricted to
effective actual rational \(D\) with
\(\operatorname{vol}(\mathcal O(D))=0\). Its prescribed full-pole mass
is zero, never an asserted \(D^n\) or volume in the big case. No
probability normalization occurs. The trivial bundle is included with
\(D=0\). Dimension zero is excluded: there the empty product has point
mass one.

\section{8. New zero-volume domain without effective
representations}\label{new-zero-volume-domain-without-effective-representations}

\textbf{Theorem 8.1 (all actual numerically trivial rational bundles).}
Suppose \(L\) is any actual rational bundle on the retained \(X\) with
\(\theta_L=0\). Retain its actual canonical metric, reference potential
zero, and the literal numerical psh definition of Section 1. Then its
entire genuine scalar cone consists of constants; these constants define
actual metrics on \(L\), not on a substitute bundle. Its independent
genuine top-degree local and global metric operator is zero, its actual
relative energy is zero, and all constants have finite energy. Its
intrinsic positive target class consists of \(\mu=0\), the equation is
actually attained and unique modulo constants, and all continuous outer
directions have zero derivative. The result covers numerical-zero actual
twists without requiring any effective divisor representation.

\subsection{8.1. The whole genuine numerical-zero
cone}\label{the-whole-genuine-numerical-zero-cone}

Let \(\phi_i\downarrow\phi\) be its one defining FS net in \(B_i\) with
\(\theta_{B_i}\to0\). Put
\(C_i=\phi_i(v_0)\downarrow C=\phi(v_0)\in\mathbb R\). Fix actual ample
\(A\). Refine by rational \(\epsilon>0\) tending to zero, with
\(\epsilon A-B_i\) actual ample. Such admissible pairs are directed and
final even if the original net is stationary: first choose a smaller
\(\epsilon\), then move into its ample tail. Multiplication by
generators and rational scaling give
\[a_{i,\epsilon}=\epsilon^{-1}(\phi_i-C_i)\in\operatorname{PSH}(A),
\qquad a_{i,\epsilon}(v_0)=0.\tag{49}\] For each divisorial \(v\), ample
compactness gives a finite evaluation lower bound \(M_v\), whence
\(\epsilon M_v\leq\phi_i(v)-C_i\leq0\). Letting \(\epsilon\to0\) proves
\(\phi(v)=C\) divisorially. Embed \(\phi\) and the constant \(C\) into
one actual ample \(H\) and use divisorial determination. Thus \(\phi=C\)
on the entire original carrier.

Constants genuinely belong to this numerical cone even when \(L\) itself
has no sections. For rational \(\epsilon\downarrow0\), the actual
\(L+\epsilon A\) is ample by numerical ampleness and has a generating
section power. Its zero-weight FS scalar is identically zero by (1).
These constants form a decreasing single net, with classes
\(\epsilon\theta_A\to0\); add \(C\). Consequently
\[\operatorname{PSH}_{\rm def}(0)=\mathbb R.\tag{50}\] This argument
explicitly distinguishes numerical approximation from a fixed-\(L\)
section net, which can be empty. The endpoint metric remains
\(e^{-C}\|\cdot\|_{\rm triv,L}\) on the original actual \(L\).

\subsection{8.2. The actual metric operator, not an assigned
zero}\label{the-actual-metric-operator-not-an-assigned-zero}

Choose an actual ample rational \(H\). Since \(\theta_L=0\), both \(H\)
and \(H+L\) are actual ample, with the same numerical class. Canonical
metric tensor functoriality gives an \textbf{actual metric identity}
\[\|\cdot\|_{\rm triv,L}=
\|\cdot\|_{\rm triv,H+L}/\|\cdot\|_{\rm triv,H}.\tag{51}\] These two
ample metrics are continuous psh-regularizable by the supplied BE
theorem. Hence (51) is a continuous dpsh metric on the actual \(L\),
with independent genuine continuous products given by BE's signed
multiadditive extension. G0 formula (5), used only on the ample factors,
computes its top product: \[\sum_{j=0}^n(-1)^{n-j}\binom nj
((H+L)^j H^{n-j})\delta_{v_0}=L^n\delta_{v_0}=0.\tag{52}\] The
intersection cancellation is literal: each numerical class
\(\theta_{H+L}\) equals \(\theta_H\). Local products are the
restrictions of these actual dpsh products, so they are zero and
invariant under changing the decomposition (51) or frames. For a tuple
of numerically trivial actual \(L_i\), express each as a difference of
two actual ample bundles; the mixed signed expansion likewise cancels to
zero. Adding constant potentials changes none of these products.

There is also an independent local \textbf{generated positive
approximation} check. On an algebraic chart choose a standard very ample
\(A\) frame so its canonical raw weight \(h_A\) is nonnegative (a
section-ratio maximum with a zero branch), and choose a frame of \(L\)
with continuous raw canonical weight \(h_L\). For positive rational
\(\epsilon\), the actual ample canonical metric of \(L+\epsilon A\) has
raw weight \[w_\epsilon=h_L+\epsilon h_A.\tag{53}\] It is an actual
finite continuous local FS weight after taking a generating ample power.
As \(\epsilon\downarrow0\), these raw weights decrease locally to
\(h_L\), uniformly on each compact frame neighborhood. Their genuine
global top products are \(\epsilon^n A^n\delta_{v_0}\) by (5), hence
tend to zero against every compact continuous local test. A finite tuple
uses \(L_i+\epsilon A_i\) and gives mass \(\epsilon^n(A_1\cdots A_n)\),
also tending to zero.

The common actual ample collar of Section 4.1 compares (53) to
\textbf{any independently supplied bounded generated local decreasing
net with the same endpoint}. Their explicit global collar endpoints
agree, so all these local limits are the same zero measure. This proves
bounded local approximation independence in exactly the generated
category, in addition to BE dpsh independence. No arbitrary
local-current density proposition is introduced. Frame changes add one
harmonic unit logarithm to the matched nets as before. Thus (52) is the
genuine original metric operator, with an actual local construction,
rather than an operator manufactured to admit a target.

\subsection{8.3. Energy, zero volume, targets and every
direction}\label{energy-zero-volume-targets-and-every-direction}

The reference is the actual canonical metric, scalar zero. The actual
relative energy of constant \(C\) is the mixed integral formula of (35)
with relative difference \(C\) and every top-degree measure zero; it
equals zero. Floors are \(\max(C,-j)\) and have the same zero products
and energy. Thus every constant is finite energy and all mixed,
decreasing and integrable quotient-weighted limits are exactly zero. The
normalized space is the singleton zero; feasibility of a scalar constant
under a continuous obstacle \(b\) means \(C\leq\inf b\), so envelopes
are actual constants and are continuous. Every original continuous outer
variation has feasible constants and zero outer energy; its two-sided
derivative is the integral against the already constructed zero
operator.

The volume really is zero. If a nonzero section of a cleared \(mL\) had
a nonzero effective zero divisor \(F\), then \(F\cdot H^{n-1}>0\) for an
ample \(H\): each positive-dimensional prime component has positive
ample intersection by the supplied ampleness criterion, and in dimension
one the positive divisor has positive degree. This contradicts
\(\theta_L=0\). Hence every such section has empty zero divisor and
trivializes \(mL\). On integral projective \(X\), global regular
functions are constants: a regular map to the affine line has proper
image in the projective line contained in the affine line, hence has
finite, irreducible image, a single point. Therefore \(h^0(X,mL)\) is
either zero or one for every cleared \(m\). The volume limsup
\(\limsup n!h^0(X,mL)/m^n\) is zero since \(n\geq1\). In particular a
nonzero numerical-zero element of
\(\operatorname{Pic}(X)\otimes\mathbb Q\) has \textbf{no} effective
rational divisor representation; an effective representative would have
to be zero and would trivialize the rational bundle.

The unique positive target of mass zero is \(\mu=0\). Its intrinsic dual
energy and all-constant integrability are finite; every actual constant
metric is an attained solution, and any two differ by a real constant.
Comparison is \(0\leq0\), while positive-mass domination is again
excluded. This theorem enlarges the zero-volume range to all actual
numerical-zero bundles, but says nothing about arbitrary
\textbf{nonzero} numerical pseudoeffective classes of zero volume.

\section{9. Primary applications, independent review and original
boundary}\label{primary-applications-independent-review-and-original-boundary}

\begin{longtable}[]{@{}
  >{\raggedright\arraybackslash}p{(\columnwidth - 4\tabcolsep) * \real{0.3333}}
  >{\raggedright\arraybackslash}p{(\columnwidth - 4\tabcolsep) * \real{0.3333}}
  >{\raggedright\arraybackslash}p{(\columnwidth - 4\tabcolsep) * \real{0.3333}}@{}}
\toprule\noalign{}
\begin{minipage}[b]{\linewidth}\raggedright
ID
\end{minipage} & \begin{minipage}[b]{\linewidth}\raggedright
Statement, hypotheses and exact application
\end{minipage} & \begin{minipage}[b]{\linewidth}\raggedright
Supplied primary body and version
\end{minipage} \\
\midrule\noalign{}
\endhead
\bottomrule\noalign{}
\endlastfoot
G0 & Full actual-ample theorem and its scoped already-glued contact
bridge, only on actual ample rational bundles over the retained
algebraically closed trivial characteristic-zero fields and projective
carrier & Entire \texttt{evidence/block-001.txt}; complete 68-body
dictionary \texttt{evidence/block-002.txt}, split under
\texttt{evidence/G0-dictionary/}. The standalone grant is not a
published nonample theorem. \\
E0 & Projective compact Hausdorff analytification; continuous canonical
metric/section norms; actual FS arithmetic and the literal single-net
numerical definition; finite-dimensional \(N^1\), open ample cone, psef
effective closed cone and big interior & BJ arXiv:1801.08229v4, March
2022. Literal sections at blocks 006--009 and dictionary units 007, 051.
No general 4.6/4.15/6.17 nonample closure is invoked. \\
E1 & On a complete normal variety over an algebraically closed field, an
R-Cartier divisor is ample iff its top intersection on every
positive-dimensional closed subvariety is positive; actual rational
ampleness is numerical & Fujino--Miyamoto arXiv:2101.00806v1, Theorem
1.4. Entire supplied pertinent statement and proof in block 010,
rendered L47--52@P1 and L179--303@P3--4. Applied only to our normal
projective original \(X\). \\
E2 & On a quasi-compact scheme with ample invertible sheaf, every
finite-type quasi-coherent sheaf has all sufficiently large twists
globally generated & Stacks tag 01Q3, supplied full Proposition 28.27.13
and proof in block 011. Used for generating section families and ideals
on projective \(X\). \\
E3 & Literal division of actual FS sections by an effective rational
Cartier section on a normal projective variety under the displayed
domination bound & BJ v4 Lemma 2.23 and complete proof, block 012;
historical inclusive LF1270--1289, review rendered L1588--1612@P26.
Those coordinate systems are not equated. \\
E4 & Semiample rational metrics are psh-regularizable; continuous
psh-regularizable metrics have genuine symmetric positive raw mixed
fundamental-cycle measures, raw intersection masses and uniform
continuity. Continuous dpsh differences have the signed multilinear
extension & BE arXiv:1805.01016v3, Sections 7.12--7.14 and 8.3--8.7,
with supplied signed extension before Example 8.11. Pertinent full
bodies at blocks 013--014, 038 and dictionary unit 036. Applies over
complete fields including the trivial valuation. Used only for actual
continuous semiample or dpsh metric products, never a nonample E1
equation. \\
E5 & On boundaryless analytic spaces over a complete field including
trivial valuation, unit logarithms have zero Hessian; the stated
continuous products are invariant under these harmonic terms & CLD
arXiv:1204.6277v2, Sections 8.8.5 and 18.1.2, supplied pertinent bodies
at blocks 015--016 and 039. Used for genuine actual frame isometries and
full-pole local products. \\
E6 & Regular local rings are UFDs; proper finite-type quasi-separated
morphisms have unique extensions of all valuation-ring diagrams & Stacks
0AG0 and 0BX5, full supplied statement/proof blocks 030--031; pertinent
existence and uniqueness units 01KA, 01KF, 01KZ and 0AFZ at blocks
032--035. Used for smooth local factoriality and centers. \\
\end{longtable}

The source table retains the author's literal extraction coordinates as
provenance. Independent review119 additionally retrieved and checked the
pertinent complete primary statements and proofs from BJ v4,
Fujino--Miyamoto v1, BE v3, CLD v2 and the named Stacks tags. Its
canonical report and exact premise/passage bindings accompany this
paper. These checks do not claim whole-book or recursive foundational
verification.

The harmless mathematical clarifications from review119 are
incorporated: the polar section need only be not identically zero
outside the reference divisor, and equation (42) explicitly uses
probability-normalized \(I_A\). The external transcription of BE Lemma
8.6 needs its curvature-difference minus sign; that transcription is not
reproduced or used with an erroneous sign here. The retained actual
signed-product identities already have their stated metric differences.

The certified scope comprises the two extensions and their indispensable
supporting Sections 1--8. It supplies no general nonnef
minimal-reference membership, full-volume equation, or general nonbig
nonzero numerical carrier. Bounded reference changes retain one
singularity type and raw mass \(A^n\); changing singularity type need
not preserve mass. The numerical-zero target is the zero measure, with
uniqueness modulo constants only because the entire genuine cone was
proved to be constants. A nonconstant zero-top-mass cone is not given
that uniqueness conclusion.

Compilation checks document production only. The frozen review task and
mathematical input hashes are retained in the accompanying provenance
manifest.

\end{document}
